\hypertarget{chapter-11-move-semantics-and-smart-pointers}{%
\section{CHAPTER 11: MOVE SEMANTICS AND SMART
POINTERS}\label{chapter-11-move-semantics-and-smart-pointers}}

\hypertarget{move-semantics-smart-pointers}{%
\section{MOVE SEMANTICS \& SMART
POINTERS}\label{move-semantics-smart-pointers}}

\hypertarget{rvalue-references-move-semantics}{%
\subsection{1. Rvalue References \& Move
Semantics}\label{rvalue-references-move-semantics}}

C++11 solves the problem of unnecessary copying with \textbf{Move
Semantics}.

\hypertarget{lvalues-vs-rvalues}{%
\subsubsection{1.1 Lvalues vs Rvalues}\label{lvalues-vs-rvalues}}

\begin{itemize}
\tightlist
\item
  \textbf{Lvalue (Left Value):} Has a name, has an address (identity).
  Persists beyond the expression.
\item
  \textbf{Rvalue (Right Value):} Temporary, no name (or about to die).
  Cannot take address.
\end{itemize}

\begin{lstlisting}[language={C++}]
int x = 10; // x is lvalue, 10 is rvalue
int y = x;  // y is lvalue
int z = x + y; // (x+y) is rvalue
\end{lstlisting}

\hypertarget{rvalue-references-t}{%
\subsubsection{\texorpdfstring{1.2 Rvalue References
(\texttt{T\&\&})}{1.2 Rvalue References (T\&\&)}}\label{rvalue-references-t}}

A reference that binds \emph{only} to rvalues. Represents an object we
can ``steal'' from.

\begin{lstlisting}[language={C++}]
void f(int& x)  { cout << "Lvalue ref"; }
void f(int&& x) { cout << "Rvalue ref"; }

int a = 5;
f(a); // Lvalue ref
f(5); // Rvalue ref
\end{lstlisting}

\hypertarget{move-constructor-move-assignment}{%
\subsubsection{1.3 Move Constructor \& Move
Assignment}\label{move-constructor-move-assignment}}

Instead of copying data (slow), we steal pointers (fast).

\begin{lstlisting}[language={C++}]
class Buffer {
    int* data;
    size_t size;
public:
    // Move Constructor
    Buffer(Buffer&& other) noexcept : data(other.data), size(other.size) {
        other.data = nullptr; // Nullify source to prevent double free
        other.size = 0;
    }

    // Move Assignment
    Buffer& operator=(Buffer&& other) noexcept {
        if (this != &other) {
            delete[] data;       // Free own resources
            data = other.data;   // Steal resources
            size = other.size;
            other.data = nullptr;// Nullify source
            other.size = 0;
        }
        return *this;
    }
};
\end{lstlisting}

\hypertarget{stdmove}{%
\subsubsection{\texorpdfstring{1.4
\texttt{std::move}}{1.4 std::move}}\label{stdmove}}

Casts an lvalue to an rvalue, enabling move semantics.

\begin{lstlisting}[language={C++}]
Buffer b1;
Buffer b2 = std::move(b1); // Calls Move Constructor
// b1 is now in valid but unspecified state (empty)
\end{lstlisting}

\textbf{Warning:} Do not use \passthrough{\lstinline!b1!} after moving
from it.

\begin{center}\rule{0.5\linewidth}{0.5pt}\end{center}

\hypertarget{smart-pointers-raii}{%
\subsection{2. Smart Pointers (RAII)}\label{smart-pointers-raii}}

Manual \passthrough{\lstinline!new!}/\passthrough{\lstinline!delete!} is
prone to leaks. C++11 introduces strict ownership models.

\hypertarget{stdunique_ptr}{%
\subsubsection{\texorpdfstring{2.1
\texttt{std::unique\_ptr}}{2.1 std::unique\_ptr}}\label{stdunique_ptr}}

\begin{itemize}
\tightlist
\item
  \textbf{Ownership:} Exclusive. Only one pointer owns the object.
\item
  \textbf{Copying:} Disabled (\passthrough{\lstinline!delete!}).
\item
  \textbf{Moving:} Allowed.
\item
  \textbf{Overhead:} Zero (same as raw pointer).
\end{itemize}

\begin{lstlisting}[language={C++}]
#include <memory>

std::unique_ptr<int> p1(new int(5));
// std::unique_ptr<int> p2 = p1; // ERROR: Cannot copy
std::unique_ptr<int> p3 = std::move(p1); // OK: Ownership transferred
\end{lstlisting}

\hypertarget{stdshared_ptr}{%
\subsubsection{\texorpdfstring{2.2
\texttt{std::shared\_ptr}}{2.2 std::shared\_ptr}}\label{stdshared_ptr}}

\begin{itemize}
\tightlist
\item
  \textbf{Ownership:} Shared (Reference Counted).
\item
  \textbf{Destruction:} Object deleted when last
  \passthrough{\lstinline!shared\_ptr!} dies.
\item
  \textbf{Overhead:} Control block on heap (ref counts).
\end{itemize}

\begin{lstlisting}[language={C++}]
auto sp1 = std::make_shared<int>(10); // Efficient allocation
auto sp2 = sp1; // Ref count = 2
\end{lstlisting}

\begin{longtable}[]{@{}
  >{\raggedright\arraybackslash}p{(\columnwidth - 0\tabcolsep) * \real{0.06}}@{}}
\toprule
\endhead
\#\#\# Professional Notes: Memory \& Resource Management \\
\#\#\#\# 1. RAII: Resource Acquisition Is Initialization RAII is the
core of modern C++ resource management. It ensures that resources
(memory, file handles, sockets) are tied to object lifetime. *
\textbf{Acquisition}: In the constructor. * \textbf{Release}: In the
destructor. * \textbf{Benefit}: Exception safety. If an exception
occurs, stack unwinding automatically calls destructors, preventing
leaks. \\
\#\#\#\# 2. The C++11 Memory Model and Atomics Before C++11, the
language had no formal definition of threads or shared memory. *
\textbf{The Problem}: Compilers and CPUs often reorder instructions for
performance, which can break multi-threaded logic. * \textbf{The
Solution}: \passthrough{\lstinline!std::atomic<T>!} and memory barriers.
They ensure that operations are visible across threads in a predictable
order (\passthrough{\lstinline!std::memory\_order!}). * \textbf{Data
Races}: Accessing the same memory location from multiple threads where
at least one is a write is \textbf{Undefined Behavior} unless
synchronized. \\
\#\#\#\# 3. Smart Pointer Gotchas *
\textbf{\passthrough{\lstinline!std::make\_shared!}
vs.~\passthrough{\lstinline!new!}}:
\passthrough{\lstinline!make\_shared!} is faster and more exception-safe
because it performs a single allocation for both the object and the
control block. * \textbf{Circular References}: If two
\passthrough{\lstinline!shared\_ptr!} objects point to each other, they
will never be deleted. Use \passthrough{\lstinline!std::weak\_ptr!} for
one of the links to break the cycle. *
\textbf{\passthrough{\lstinline!this!} as shared}: Use
\passthrough{\lstinline!std::enable\_shared\_from\_this<T>!} if you need
to pass a \passthrough{\lstinline!shared\_ptr!} to the current object
from inside a member function. \\
\bottomrule
\end{longtable}

\hypertarget{stdweak_ptr}{%
\subsubsection{\texorpdfstring{2.3
\texttt{std::weak\_ptr}}{2.3 std::weak\_ptr}}\label{stdweak_ptr}}

\begin{itemize}
\tightlist
\item
  \textbf{Ownership:} Non-owning observer of
  \passthrough{\lstinline!shared\_ptr!}.
\item
  \textbf{Use Case:} Break cyclic references (A-\textgreater B,
  B-\textgreater A).
\end{itemize}

\begin{lstlisting}[language={C++}]
std::shared_ptr<int> sp = std::make_shared<int>(42);
std::weak_ptr<int> wp = sp;

if (auto locked = wp.lock()) { // Check if object exists
    std::cout << *locked;
}
\end{lstlisting}

\hypertarget{stdmake_shared-vs-new}{%
\subsubsection{\texorpdfstring{2.4 \texttt{std::make\_shared} vs
\texttt{new}}{2.4 std::make\_shared vs new}}\label{stdmake_shared-vs-new}}

Always use \passthrough{\lstinline!std::make\_shared!}. It allocates the
object and control block in a single heap allocation (cache efficient).

\begin{center}\rule{0.5\linewidth}{0.5pt}\end{center}

\hypertarget{perfect-forwarding}{%
\subsection{3. Perfect Forwarding}\label{perfect-forwarding}}

Used in templates to preserve value category (lvalue vs rvalue).

\begin{lstlisting}[language={C++}]
template<typename T>
void wrapper(T&& arg) {
    // std::forward passes arg as lvalue if given lvalue,
    // rvalue if given rvalue.
    func(std::forward<T>(arg));
}
\end{lstlisting}

\hypertarget{professional-notes-chapter-33-smart-pointers}{%
\section{Professional Notes: Chapter 33: Smart
Pointers}\label{professional-notes-chapter-33-smart-pointers}}

Section 33.1: Unique ownership (std::unique\_ptr) Section 33.2: Sharing
ownership (std::shared\_ptr) Section 33.3: Sharing with temporary
ownership (std::weak\_ptr) Section 33.4: Using custom deleters to create
a wrapper to a C interface Section 33.5: Unique ownership without move
semantics (auto\_ptr) Section 33.6: Casting std::shared\_ptr pointers
Section 33.7: Writing a smart pointer: value\_ptr Section 33.8: Getting
a shared\_ptr referring to this

\hypertarget{professional-notes-chapter-106-move-semantics}{%
\section{Professional Notes: Chapter 106: Move
Semantics}\label{professional-notes-chapter-106-move-semantics}}

Section 106.1: Move semantics Section 106.2: Using std::move to reduce
complexity from O(n) to O(n) Section 106.3: Move constructor Section
106.4: Re-use a moved object Section 106.5: Move assignment Section
106.6: Using move semantics on containers

\hypertarget{professional-notes-chapter-33-smart-pointers-1}{%
\section{Professional Notes: Chapter 33: Smart
Pointers}\label{professional-notes-chapter-33-smart-pointers-1}}

Section 33.1: Unique ownership (std::unique\_ptr) Version C++11 A
std::unique\_ptr is a class template that manages the lifetime of a
dynamically stored object. Unlike for std::shared\_ptr, the dynamic
object is owned by only one instance of a std::unique\_ptr at any time,
// Creates a dynamic int with value of 20 owned by a unique pointer
std::unique\_ptr ptr = std::make\_unique(20); (Note: std::unique\_ptr is
available since C++11 and std::make\_unique since C++14.) Only the
variable ptr holds a pointer to a dynamically allocated int. When a
unique pointer that owns an object goes out of scope, the owned object
is deleted, i.e.~its destructor is called if the object is of class
type, and the memory for that object is released. To use
std::unique\_ptr and std::make\_unique with array-types, use their array
specializations: // Creates a unique\_ptr to an int with value 59
std::unique\_ptr ptr = std::make\_unique(59); // Creates a unique\_ptr
to an array of 15 ints std::unique\_ptr\textless int{[}{]}\textgreater{}
ptr = std::make\_unique\textless int{[}{]}\textgreater(15); You can
access the std::unique\_ptr just like a raw pointer, because it
overloads those operators. You can transfer ownership of the contents of
a smart pointer to another pointer by using std::move, which will cause
the original smart pointer to point to nullptr. // 1. std::unique\_ptr
std::unique\_ptr ptr = std::make\_unique(); // Change value to 1
\emph{ptr = 1; // 2. std::unique\_ptr (by moving `ptr' to `ptr2', `ptr'
doesn't own the object anymore) std::unique\_ptr ptr2 = std::move(ptr);
int a = }ptr2; // `a' is 1 int b = \emph{ptr; // undefined behavior!
`ptr' is `nullptr' // (because of the move command above) Passing
unique\_ptr to functions as parameter: void foo(std::unique\_ptr ptr) \{
// Your code goes here \} std::unique\_ptr ptr = std::make\_unique(59);
foo(std::move(ptr)) Returning unique\_ptr from functions. This is the
preferred C++11 way of writing factory functions, as it clearly conveys
the ownership semantics of the return: the caller owns the resulting
unique\_ptr and is responsible for it. std::unique\_ptr foo() \{
std::unique\_ptr ptr = std::make\_unique(59); return ptr; \}
std::unique\_ptr ptr = foo(); Compare this to: int} foo\_cpp03(); int* p
= foo\_cpp03(); // do I own p? do I have to delete it at some point? //
it's not readily apparent what the answer is. Version \textless{} C++14
The class template make\_unique is provided since C++14. It's easy to
add it manually to C++11 code: template\textless typename T,
typename\ldots{} Args\textgreater{} typename
std::enable\_if\textless!std::is\_array::value,
std::unique\_ptr\textgreater::type make\_unique(Args\&\&\ldots{} args)
\{ return std::unique\_ptr(new T(std::forward(args)\ldots)); \} // Use
make\_unique for arrays template typename
std::enable\_if\textless std::is\_array::value,
std::unique\_ptr\textgreater::type make\_unique(size\_t n) \{ return
std::unique\_ptr(new typename std::remove\_extent::type\href{}{n}); \}
Version C++11 Unlike the dumb smart pointer (std::auto\_ptr),
unique\_ptr can also be instantiated with vector allocation (not
std::vector). Earlier examples were for scalar allocations. For example
to have a dynamically allocated integer array for 10 elements, you would
specify int{[}{]} as the template type (and not just int):
std::unique\_ptr\textless int{[}{]}\textgreater{} arr\_ptr =
std::make\_unique\textless int{[}{]}\textgreater(10); Which can be
simplied with: auto arr\_ptr =
std::make\_unique\textless int{[}{]}\textgreater(10); Now, you use
arr\_ptr as if it is an array: arr\_ptr{[}2{]} = 10; // Modify third
element You need not to worry about de-allocation. This template
specialized version calls constructors and destructors appropriately.
Using vectored version of unique\_ptr or a vector itself - is a personal
choice. In versions prior to C++11, std::auto\_ptr was available. Unlike
unique\_ptr it is allowed to copy auto\_ptrs, upon which the source ptr
will lose the ownership of the contained pointer and the target receives
it. Section 33.2: Sharing ownership (std::shared\_ptr) The class
template std::shared\_ptr denes a shared pointer that is able to share
ownership of an object with other shared pointers. This contrasts to
std::unique\_ptr which represents exclusive ownership. The sharing
behavior is implemented through a technique known as reference counting,
where the number of shared pointers that point to the object is stored
alongside it. When this count reaches zero, either through the
destruction or reassignment of the last std::shared\_ptr instance, the
object is automatically destroyed. // Creation: `firstShared' is a
shared pointer for a new instance of `Foo' std::shared\_ptr firstShared
= std::make\_shared(/\emph{args}/); To create multiple smart pointers
that share the same object, we need to create another shared\_ptr that
aliases the rst shared pointer. Here are 2 ways of doing it:
std::shared\_ptr secondShared(firstShared); // 1st way: Copy
constructing std::shared\_ptr secondShared; secondShared = firstShared;
// 2nd way: Assigning Either of the above ways makes secondShared a
shared pointer that shares ownership of our instance of Foo with
firstShared. The smart pointer works just like a raw pointer. This
means, you can use * to dereference them. The regular -\textgreater{}
operator works as well: secondShared-\textgreater test(); // Calls
Foo::test() Finally, when the last aliased shared\_ptr goes out of
scope, the destructor of our Foo instance is called. Warning:
Constructing a shared\_ptr might throw a bad\_alloc exception when extra
data for shared ownership semantics needs to be allocated. If the
constructor is passed a regular pointer it assumes to own the object
pointed to and calls the deleter if an exception is thrown. This means
shared\_ptr(new T(args)) will not leak a T object if allocation of
shared\_ptr fails. However, it is advisable to use make\_shared(args) or
allocate\_shared(alloc, args), which enable the implementation to
optimize the memory allocation. Allocating Arrays({[}{]}) using
shared\_ptr Version C++11 Version \textless{} C++17 Unfortunately, there
is no direct way to allocate Arrays using
make\_shared\textless\textgreater. It is possible to create arrays for
shared\_ptr\textless\textgreater{} using new and std::default\_delete.
For example, to allocate an array of 10 integers, we can write the code
as shared\_ptr sh(new int{[}10{]},
std::default\_delete\textless int{[}{]}\textgreater()); Specifying
std::default\_delete is mandatory here to make sure that the allocated
memory is correctly cleaned up using delete{[}{]}. If we know the size
at compile time, we can do it this way: template struct
shared\_array\_maker \{\}; template\textless class T, std::size\_t
N\textgreater{} struct
shared\_array\_maker\textless T{[}N{]}\textgreater{} \{ std::shared\_ptr
operator()const\{  auto r =
std::make\_shared\textless std::array\textless T,N\textgreater\textgreater();
if (!r) return \{\}; return \{r.data(), r\}; \} \}; template auto
make\_shared\_array() -\textgreater{} decltype(
shared\_array\_maker\{\}() ) \{ return shared\_array\_maker\{\}(); \}
then make\_shared\_array\textless int{[}10{]}\textgreater{} returns a
shared\_ptr pointing to 10 ints all default constructed. Version C++17
With C++17, shared\_ptr gained special support for array types. It is no
longer necessary to specify the array-deleter explicitly, and the shared
pointer can be dereferenced using the {[}{]} array index operator:
std::shared\_ptr\textless int{[}{]}\textgreater{} sh(new int{[}10{]});
sh{[}0{]} = 42; Shared pointers can point to a sub-object of the object
it owns: struct Foo \{ int x; \}; std::shared\_ptr p1 =
std::make\_shared(); std::shared\_ptr p2(p1, \&p1-\textgreater x); Both
p2 and p1 own the object of type Foo, but p2 points to its int member x.
This means that if p1 goes out of scope or is reassigned, the underlying
Foo object will still be alive, ensuring that p2 does not dangle.
Important: A shared\_ptr only knows about itself and all other
shared\_ptr that were created with the alias constructor. It does not
know about any other pointers, including all other shared\_ptrs created
with a reference to the same Foo instance: Foo \emph{foo = new Foo;
std::shared\_ptr shared1(foo); std::shared\_ptr shared2(foo); // don't
do this shared1.reset(); // this will delete foo, since shared1 // was
the only shared\_ptr that owned it shared2-\textgreater test(); //
UNDEFINED BEHAVIOR: shared2's foo has been // deleted already!!
Ownership Transfer of shared\_ptr By default, shared\_ptr increments the
reference count and doesn't transfer the ownership. However, it can be
made to transfer the ownership using std::move: shared\_ptr up =
make\_shared(); // Transferring the ownership shared\_ptr up2 =
move(up); // At this point, the reference count of up = 0 and the //
ownership of the pointer is solely with up2 with reference count = 1
Section 33.3: Sharing with temporary ownership (std::weak\_ptr)
Instances of std::weak\_ptr can point to objects owned by instances of
std::shared\_ptr while only becoming temporary owners themselves. This
means that weak pointers do not alter the object's reference count and
therefore do not prevent an object's deletion if all of the object's
shared pointers are reassigned or destroyed. In the following example
instances of std::weak\_ptr are used so that the destruction of a tree
object is not inhibited: \#include \#include struct TreeNode \{
std::weak\_ptr parent; std::vector\textless{} std::shared\_ptr
\textgreater{} children; \}; int main() \{ // Create a TreeNode to serve
as the root/parent. std::shared\_ptr root(new TreeNode); // Give the
parent 100 child nodes. for (size\_t i = 0; i \textless{} 100; ++i) \{
std::shared\_ptr child(new TreeNode);
root-\textgreater children.push\_back(child); child-\textgreater parent
= root; \} // Reset the root shared pointer, destroying the root object,
and // subsequently its child nodes. root.reset(); \} As child nodes are
added to the root node's children, their std::weak\_ptr member parent is
set to the root node. The member parent is declared as a weak pointer as
opposed to a shared pointer such that the root node's reference count is
not incremented. When the root node is reset at the end of main(), the
root is destroyed. Since the only remaining std::shared\_ptr references
to the child nodes were contained in the root's collection children, all
child nodes are subsequently destroyed as well. Due to control block
implementation details, shared\_ptr allocated memory may not be released
until shared\_ptr reference counter and weak\_ptr reference counter both
reach zero. \#include int main() \{ \{ std::weak\_ptr wk; \{ //
std::make\_shared is optimized by allocating only once // while
std::shared\_ptr(new int(42)) allocates twice. // Drawback of
std::make\_shared is that control block is tied to our integer
std::shared\_ptr sh = std::make\_shared(42); wk = sh; // sh memory
should be released at this point\ldots{} \} // \ldots{} but wk is still
alive and needs access to control block \} // now memory is released (sh
and wk) \} Since std::weak\_ptr does not keep its referenced object
alive, direct data access through a std::weak\_ptr is not possible.
Instead it provides a lock() member function that attempts to retrieve a
std::shared\_ptr to the referenced object: \#include \#include int
main() \{ \{ std::weak\_ptr wk; std::shared\_ptr sp; \{ std::shared\_ptr
sh = std::make\_shared(42); wk = sh; // calling lock will create a
shared\_ptr to the object referenced by wk sp = wk.lock(); // sh will be
destroyed after this point, but sp is still alive \} // sp still keeps
the data alive. // At this point we could even call lock() again // to
retrieve another shared\_ptr to the same data from wk assert(}sp == 42);
assert(!wk.expired()); // resetting sp will delete the data, // as it is
currently the last shared\_ptr with ownership sp.reset(); // attempting
to lock wk now will return an empty shared\_ptr, // as the data has
already been deleted sp = wk.lock(); assert(!sp); assert(wk.expired());
\} \} Section 33.4: Using custom deleters to create a wrapper to a C
interface Many C interfaces such as SDL2 have their own deletion
functions. This means that you cannot use smart pointers directly:
std::unique\_ptr a; // won't work, UNSAFE! Instead, you need to dene
your own deleter. The examples here use the SDL\_Surface structure which
should be freed using the SDL\_FreeSurface() function, but they should
be adaptable to many other C interfaces. The deleter must be callable
with a pointer argument, and therefore can be e.g.~a simple function
pointer: std::unique\_ptr\textless SDL\_Surface,
void(\emph{)(SDL\_Surface})\textgreater{} a(pointer, SDL\_FreeSurface);
Any other callable object will work, too, for example a class with an
operator(): struct SurfaceDeleter \{ void operator()(SDL\_Surface* surf)
\{ SDL\_FreeSurface(surf); \} \};
std::unique\_ptr\textless SDL\_Surface, SurfaceDeleter\textgreater{}
a(pointer, SurfaceDeleter\{\}); // safe
std::unique\_ptr\textless SDL\_Surface, SurfaceDeleter\textgreater{}
b(pointer); // equivalent to the above // as the deleter is
value-initialized This not only provides you with safe, zero overhead
(if you use unique\_ptr) automatic memory management, you also get
exception safety. Note that the deleter is part of the type for
unique\_ptr, and the implementation can use the empty base optimization
to avoid any change in size for empty custom deleters. So while
std::unique\_ptr\textless SDL\_Surface, SurfaceDeleter\textgreater{} and
std::unique\_ptr\textless SDL\_Surface,
void(\emph{)(SDL\_Surface})\textgreater{} solve the same problem in a
similar way, the former type is still only the size of a pointer while
the latter type has to hold two pointers: both the SDL\_Surface* and the
function pointer! When having free function custom deleters, it is
preferable to wrap the function in an empty type. In cases where
reference counting is important, one could use a shared\_ptr instead of
an unique\_ptr. The shared\_ptr always stores a deleter, this erases the
type of the deleter, which might be useful in APIs. The disadvantages of
using shared\_ptr over unique\_ptr include a higher memory cost for
storing the deleter and a performance cost for maintaining the reference
count. // deleter required at construction time and is part of the type
std::unique\_ptr\textless SDL\_Surface,
void(\emph{)(SDL\_Surface})\textgreater{} a(pointer, SDL\_FreeSurface);
// deleter is only required at construction time, not part of the type
std::shared\_ptr b(pointer, SDL\_FreeSurface); Version C++17 With
template auto, we can make it even easier to wrap our custom deleters:
template struct FunctionDeleter \{ template void operator()(T* ptr) \{
DeleteFn(ptr); \} \}; template \textless class T, auto
DeleteFn\textgreater{} using unique\_ptr\_deleter =
std::unique\_ptr\textless T, FunctionDeleter\textgreater; With which the
above example is simply: unique\_ptr\_deleter\textless SDL\_Surface,
SDL\_FreeSurface\textgreater{} c(pointer); Here, the purpose of auto is
to handle all free functions, whether they return void
(e.g.~SDL\_FreeSurface) or not (e.g.~fclose). Section 33.5: Unique
ownership without move semantics (auto\_ptr) Version \textless{} C++11
NOTE: std::auto\_ptr has been deprecated in C++11 and will be removed in
C++17. You should only use this if you are forced to use C++03 or
earlier and are willing to be careful. It is recommended to move to
unique\_ptr in combination with std::move to replace std::auto\_ptr
behavior. Before we had std::unique\_ptr, before we had move semantics,
we had std::auto\_ptr. std::auto\_ptr provides unique ownership but
transfers ownership upon copy. As with all smart pointers,
std::auto\_ptr automatically cleans up resources (see RAII): \{
std::auto\_ptr p(new int(42)); std::cout \textless\textless{} \emph{p;
\} // p is deleted here, no memory leaked but allows only one owner:
std::auto\_ptr px = \ldots; std::auto\_ptr py = px; // px is now empty
This allows to use std::auto\_ptr to keep ownership explicit and unique
at the danger of losing ownership unintended: void f(std::auto\_ptr ) \{
// assumes ownership of X // deletes it at end of scope \};
std::auto\_ptr px = \ldots; f(px); // f acquires ownership of underlying
X // px is now empty px-\textgreater foo(); // NPE! //
px.\textasciitilde auto\_ptr() does NOT delete The transfer of ownership
happened in the ``copy'' constructor. auto\_ptr's copy constructor and
copy assignment operator take their operands by non-const reference so
that they could be modied. An example implementation might be: template
class auto\_ptr \{ T} ptr; public: auto\_ptr(auto\_ptr\& rhs) :
ptr(rhs.release()) \{ \} auto\_ptr\& operator=(auto\_ptr\& rhs) \{
reset(rhs.release()); return \emph{this; \} T} release() \{ T* tmp =
ptr; ptr = nullptr; return tmp; \} void reset(T* tmp = nullptr) \{ if
(ptr != tmp) \{ delete ptr; ptr = tmp; \} \}  /* other functions
\ldots{} \emph{/ \}; This breaks copy semantics, which require that
copying an object leaves you with two equivalent versions of it. For any
copyable type, T, I should be able to write: T a = \ldots; T b(a);
assert(b == a); But for auto\_ptr, this is not the case. As a result, it
is not safe to put auto\_ptrs in containers. Section 33.6: Casting
std::shared\_ptr pointers It is not possible to directly use
static\_cast, const\_cast, dynamic\_cast and reinterpret\_cast on
std::shared\_ptr to retrieve a pointer sharing ownership with the
pointer being passed as argument. Instead, the functions
std::static\_pointer\_cast, std::const\_pointer\_cast,
std::dynamic\_pointer\_cast and std::reinterpret\_pointer\_cast should
be used: struct Base \{ virtual \textasciitilde Base() noexcept \{\};
\}; struct Derived: Base \{\}; auto derivedPtr(std::make\_shared());
auto basePtr(std::static\_pointer\_cast(derivedPtr)); auto
constBasePtr(std::const\_pointer\_cast(basePtr)); auto
constDerivedPtr(std::dynamic\_pointer\_cast(constBasePtr)); Note that
std::reinterpret\_pointer\_cast is not available in C++11 and C++14, as
it was only proposed by N3920 and adopted into Library Fundamentals TS
in February 2014. However, it can be implemented as follows: template
\textless typename To, typename From\textgreater{} inline
std::shared\_ptr reinterpret\_pointer\_cast( std::shared\_ptr const \&
ptr) noexcept \{ return std::shared\_ptr(ptr,
reinterpret\_cast\textless To }\textgreater(ptr.get())); \} Section
33.7: Writing a smart pointer: value\_ptr A value\_ptr is a smart
pointer that behaves like a value. When copied, it copies its contents.
When created, it creates its contents. // Like std::default\_delete:
template struct default\_copier \{ // a copier must handle a null T
const* in and return null: T* operator()(T const* tin)const \{ if (!tin)
return nullptr; return new T(\emph{tin); \} void operator()(void} dest,
T const* tin)const \{ if (!tin) return; return new(dest) T(\emph{tin);
\} \}; // tag class to handle empty case: struct empty\_ptr\_t \{\};
constexpr empty\_ptr\_t empty\_ptr\{\}; // the value pointer type
itself: template\textless class T, class Copier=default\_copier, class
Deleter=std::default\_delete, class Base=std::unique\_ptr\textless T,
Deleter\textgreater{} \textgreater{} struct value\_ptr:Base, private
Copier \{ using copier\_type=Copier; // also typedefs from unique\_ptr
using Base::Base; value\_ptr( T const\& t ): Base( std::make\_unique(t)
), Copier() \{\} value\_ptr( T \&\& t ): Base(
std::make\_unique(std::move(t)) ), Copier() \{\} // almost-never-empty:
value\_ptr(): Base( std::make\_unique() ), Copier() \{\} value\_ptr(
empty\_ptr\_t ) \{\} value\_ptr( Base b, Copier c=\{\} ):
Base(std::move(b)), Copier(std::move(c)) \{\} Copier const\&
get\_copier() const \{ return }this; \} value\_ptr clone() const \{
return \{ Base( get\_copier()(this-\textgreater get()),
this-\textgreater get\_deleter() ), get\_copier() \}; \}
value\_ptr(value\_ptr\&\&)=default; value\_ptr\&
operator=(value\_ptr\&\&)=default; value\_ptr(value\_ptr const\&
o):value\_ptr(o.clone()) \{\} value\_ptr\& operator=(value\_ptr
const\&o) \{ if (o \&\& \emph{this) \{ // if we are both non-null,
assign contents: \textbf{this = \emph{o; \} else \{ // otherwise, assign
a clone (which could itself be null): }this = o.clone(); \} return
\emph{this; \} value\_ptr\& operator=( T const\& t ) \{ if (}this) \{
}this = t; \} else \{ }this = value\_ptr(t); \}  return \emph{this; \}
value\_ptr\& operator=( T \&\& t ) \{ if (}this) \{ \textbf{this =
std::move(t); \} else \{ \emph{this = value\_ptr(std::move(t)); \}
return }this; \} T\& get() \{ return }this; \} T const\& get() const \{
return **this; \} T* get\_pointer() \{ if (!\emph{this) return nullptr;
return std::addressof(get()); \} T const} get\_pointer() const \{ if
(!\emph{this) return nullptr; return std::addressof(get()); \} //
operator-\textgreater{} from unique\_ptr \}; template\textless class T,
class\ldots Args\textgreater{} value\_ptr make\_value\_ptr(
Args\&\&\ldots{} args ) \{ return
\{std::make\_unique(std::forward(args)\ldots)\}; \} This particular
value\_ptr is only empty if you construct it with empty\_ptr\_t or if
you move from it. It exposes the fact it is a unique\_ptr, so explicit
operator bool() const works on it. .get() has been changed to return a
reference (as it is almost never empty), and .get\_pointer() returns a
pointer instead. This smart pointer can be useful for pImpl cases, where
we want value-semantics but we also don't want to expose the contents of
the pImpl outside of the implementation le. With a non-default Copier,
it can even handle virtual base classes that know how to produce
instances of their derived and turn them into value-types. Section 33.8:
Getting a shared\_ptr referring to this enable\_shared\_from\_this
enables you to get a valid shared\_ptr instance to this. By deriving
your class from the class template enable\_shared\_from\_this, you
inherit a method shared\_from\_this that returns a shared\_ptr instance
to this. Note that the object must be created as a shared\_ptr in rst
place: \#include class A: public enable\_shared\_from\_this \{ \}; A}
ap1 =new A(); shared\_ptr ap2(ap1); // First prepare a shared pointer to
the object and hold it! // Then get a shared pointer to the object from
the object itself shared\_ptr ap3 =
ap1-\textgreater shared\_from\_this(); int c3 =ap3.use\_count(); // =2:
pointing to the same object Note(2) you cannot call
enable\_shared\_from\_this inside the constructor. \#include //
enable\_shared\_from\_this class Widget : public
std::enable\_shared\_from\_this\textless{} Widget \textgreater{} \{
public: void DoSomething() \{ std::shared\_ptr\textless{} Widget
\textgreater{} self = shared\_from\_this(); someEvent -\textgreater{}
Register( self ); \} private: \}; int main() \{ auto w =
std::make\_shared\textless{} Widget \textgreater(); w -\textgreater{}
DoSomething(); \} If you use shared\_from\_this() on an object not owned
by a shared\_ptr, such as a local automatic object or a global object,
then the behavior is undened. Since C++17 it throws std::bad\_alloc
instead. Using shared\_from\_this() from a constructor is equivalent to
using it on an object not owned by a shared\_ptr, because the objects is
possessed by the shared\_ptr after the constructor returns.

\hypertarget{professional-notes-chapter-106-move-semantics-1}{%
\section{Professional Notes: Chapter 106: Move
Semantics}\label{professional-notes-chapter-106-move-semantics-1}}

Section 106.1: Move semantics Move semantics are a way of moving one
object to another in C++. For this, we empty the old object and place
everything it had in the new object. For this, we must understand what
an rvalue reference is. An rvalue reference (T\&\& where T is the object
type) is not much dierent than a normal reference (T\&, now called
lvalue references). But they act as 2 dierent types, and so, we can make
constructors or functions that take one type or the other, which will be
necessary when dealing with move semantics. The reason why we need two
dierent types is to specify two dierent behaviors. Lvalue reference
constructors are related to copying, while rvalue reference constructors
are related to moving. To move an object, we will use std::move(obj).
This function returns an rvalue reference to the object, so that we can
steal the data from that object into a new one. There are several ways
of doing this which are discussed below. Important to note is that the
use of std::move creates just an rvalue reference. In other words the
statement std::move(obj) does not change the content of obj, while auto
obj2 = std::move(obj) (possibly) does. Section 106.2: Using std::move to
reduce complexity from O(n) to O(n) C++11 introduced core language and
standard library support for moving an object. The idea is that when an
object o is a temporary and one wants a logical copy, then its safe to
just pilfer o's resources, such as a dynamically allocated buer, leaving
o logically empty but still destructible and copyable. The core language
support is mainly the rvalue reference type builder \&\&, e.g.,
std::string\&\& is an rvalue reference to a std::string, indicating that
that referred to object is a temporary whose resources can just be
pilfered (i.e.~moved) special support for a move constructor T( T\&\& ),
which is supposed to eciently move resources from the specied other
object, instead of actually copying those resources, and special support
for a move assignment operator auto operator=(T\&\&) -\textgreater{}
T\&, which also is supposed to move from the source. The standard
library support is mainly the std::move function template from the
header. This function produces an rvalue reference to the specied
object, indicating that it can be moved from, just as if it were a
temporary. For a container actual copying is typically of O(n)
complexity, where n is the number of items in the container, while
moving is O(1), constant time. And for an algorithm that logically
copies that container n times, this can reduce the complexity from the
usually impractical O(n) to just linear O(n). In his article Containers
That Never Change in Dr.~Dobbs Journal in September 19 2013, Andrew
Koenig presented an interesting example of algorithmic ineciency when
using a style of programming where variables are immutable after
initialization. With this style loops are generally expressed using
recursion. And for some algorithms such as generating a Collatz
sequence, the recursion requires logically copying a container: // Based
on an example by Andrew Koenig in his Dr.~Dobbs Journal article //
Containers That Never Change September 19, 2013, available at //
\textless url:
http://www.drdobbs.com/cpp/containters-that-never-change/240161543\textgreater{}
// Includes here, e.g.~ namespace my \{ template\textless{} class Item
\textgreater{} using Vector\_ = /* E.g. std::vector \emph{/; auto
concat( Vector\_ const\& v, int const x ) -\textgreater{} Vector\_ \{
auto result\{ v \}; result.push\_back( x ); return result; \} auto
collatz\_aux( int const n, Vector\_ const\& result ) -\textgreater{}
Vector\_ \{ if( n == 1 ) \{ return result; \} auto const new\_result =
concat( result, n ); if( n \% 2 == 0 ) \{ return collatz\_aux( n/2,
new\_result ); \} else \{ return collatz\_aux( 3}n + 1, new\_result );
\} \} auto collatz( int const n ) -\textgreater{} Vector\_ \{ assert( n
!= 0 ); return collatz\_aux( n, Vector\_() ); \} \} // namespace my
\#include using namespace std; auto main() -\textgreater{} int \{ for(
int const x : my::collatz( 42 ) ) \{ cout \textless\textless{} x
\textless\textless{} ' `; \} cout \textless\textless{}'\n`; \} Output:
42 21 64 32 16 8 4 2 The number of item copy operations due to copying
of the vectors is here roughly O(n), since it's the sum 1 + 2 + 3 +
\ldots{} n. In concrete numbers, with g++ and Visual C++ compilers the
above invocation of collatz(42) resulted in a Collatz sequence of 8
items and 36 item copy operations (8\emph{7/2 = 28, plus some) in vector
copy constructor calls. All of these item copy operations can be removed
by simply moving vectors whose values are not needed anymore. To do this
it's necessary to remove const and reference for the vector type
arguments, passing the vectors by value. The function returns are
already automatically optimized. For the calls where vectors are passed,
and not used again further on in the function, just apply std::move to
move those buers rather than actually copying them: using std::move;
auto concat( Vector\_ v, int const x ) -\textgreater{} Vector\_ \{
v.push\_back( x ); // warning: moving a local object in a return
statement prevents copy elision {[}-Wpessimizing- move{]} // See
https://stackoverflow.com/documentation/c\%2b\%2b/2489/copy-elision //
return move( v ); return v; \} auto collatz\_aux( int const n, Vector\_
result ) -\textgreater{} Vector\_ \{ if( n == 1 ) \{ return result; \}
auto new\_result = concat( move( result ), n ); struct result; // Make
absolutely sure no use of \passthrough{\lstinline!result!} after this.
if( n \% 2 == 0 ) \{ return collatz\_aux( n/2, move( new\_result ) ); \}
else \{ return collatz\_aux( 3}n + 1, move( new\_result ) ); \} \} auto
collatz( int const n ) -\textgreater{} Vector\_ \{ assert( n != 0 );
return collatz\_aux( n, Vector\_() ); \} Here, with g++ and Visual C++
compilers, the number of item copy operations due to vector copy
constructor invocations, was exactly 0. The algorithm is necessarily
still O(n) in the length of the Collatz sequence produced, but this is a
quite dramatic improvement: O(n) O(n). With some language support one
could perhaps use moving and still express and enforce the immutability
of a variable between its initialization and nal move, after which any
use of that variable should be an error. Alas, as of C++14 C++ does not
support that. For loop-free code the no use after move can be enforced
via a re-declaration of the relevant name as an incomplete struct, as
with struct result; above, but this is ugly and not likely to be
understood by other programmers; also the diagnostics can be quite
misleading. Summing up, the C++ language and library support for moving
allows drastic improvements in algorithm complexity, but due the
support's incompleteness, at the cost of forsaking the code correctness
guarantees and code clarity that const can provide. For completeness,
the instrumented vector class used to measure the number of item copy
operations due to copy constructor invocations: template\textless{}
class Item \textgreater{} class Copy\_tracking\_vector \{ private:
static auto n\_copy\_ops() -\textgreater{} int\& \{ static int value;
return value; \} vector items\_; public: static auto n() -\textgreater{}
int \{ return n\_copy\_ops(); \} void push\_back( Item const\& o ) \{
items\_.push\_back( o ); \} auto begin() const \{ return
items\_.begin(); \} auto end() const \{ return items\_.end(); \}
Copy\_tracking\_vector()\{\} Copy\_tracking\_vector(
Copy\_tracking\_vector const\& other ) : items\_( other.items\_ ) \{
n\_copy\_ops() += items\_.size(); \} Copy\_tracking\_vector(
Copy\_tracking\_vector\&\& other ) : items\_( move( other.items\_ ) )
\{\} \}; Section 106.3: Move constructor Say we have this code snippet.
class A \{ public: int a; int b; A(const A \&other) \{
this-\textgreater a = other.a; this-\textgreater b = other.b; \} \}; To
create a copy constructor, that is, to make a function that copies an
object and creates a new one, we normally would choose the syntax shown
above, we would have a constructor for A that takes an reference to
another object of type A, and we would copy the object manually inside
the method. Alternatively, we could have written A(const A \&) =
default; which automatically copies over all members, making use of its
copy constructor. To create a move constructor, however, we will be
taking an rvalue reference instead of an lvalue reference, like here.
class Wallet \{ public: int nrOfDollars; Wallet() = default; //default
ctor Wallet(Wallet \&\&other) \{ this-\textgreater nrOfDollars =
other.nrOfDollars; other.nrOfDollars = 0; \} \}; Please notice that we
set the old values to zero. The default move constructor
(Wallet(Wallet\&\&) = default;) copies the value of nrOfDollars, as it
is a POD. As move semantics are designed to allow 'stealing' state from
the original instance, it is important to consider how the original
instance should look like after this stealing. In this case, if we would
not change the value to zero we would have doubled the amount of dollars
into play. Wallet a; a.nrOfDollars = 1; Wallet b (std::move(a));
//calling B(B\&\& other); std::cout \textless\textless{} a.nrOfDollars
\textless\textless{} std::endl; //0 std::cout \textless\textless{}
b.nrOfDollars \textless\textless{} std::endl; //1 Thus we have move
constructed an object from an old one. While the above is a simple
example, it shows what the move constructor is intended to do. It
becomes more useful in more complex cases, such as when resource
management is involved. // Manages operations involving a specified
type. // Owns a helper on the heap, and one in its memory (presumably on
the stack). // Both helpers are DefaultConstructible, CopyConstructible,
and MoveConstructible. template\textless typename T, template typename
HeapHelper, template typename StackHelper\textgreater{} class
OperationsManager \{ using MyType = OperationsManager\textless T,
HeapHelper, StackHelper\textgreater; HeapHelper* h\_helper; StackHelper
s\_helper; // \ldots{} public: // Default constructor \& Rule of Five.
OperationsManager() : h\_helper(new HeapHelper) \{\}
OperationsManager(const MyType\& other) : h\_helper(new
HeapHelper(\emph{other.h\_helper)), s\_helper(other.s\_helper) \{\}
MyType\& operator=(MyType copy) \{ swap(}this, copy); return \emph{this;
\} \textasciitilde OperationsManager() \{ if (h\_helper) \{ delete
h\_helper; \}  \} // Move constructor (without swap()). // Takes
other's HeapHelper}. // Takes other's StackHelper, by forcing the use of
StackHelper's move constructor. // Replaces other's HeapHelper* with
nullptr, to keep other from deleting our shiny // new helper when it's
destroyed. OperationsManager(MyType\&\& other) noexcept :
h\_helper(other.h\_helper), s\_helper(std::move(other.s\_helper)) \{
other.h\_helper = nullptr; \} // Move constructor (with swap()). //
Places our members in the condition we want other's to be in, then
switches members // with other. // OperationsManager(MyType\&\& other)
noexcept : h\_helper(nullptr) \{ // swap(\emph{this, other); // \} //
Copy/move helper. friend void swap(MyType\& left, MyType\& right)
noexcept \{ std::swap(left.h\_helper, right.h\_helper);
std::swap(left.s\_helper, right.s\_helper); \} \}; Section 106.4: Re-use
a moved object You can re-use a moved object: void
consumingFunction(std::vector vec) \{ // Some operations \} int main()
\{ // initialize vec with 1, 2, 3, 4 std::vector vec\{1, 2, 3, 4\}; //
Send the vector by move consumingFunction(std::move(vec)); // Here the
vec object is in an indeterminate state. // Since the object is not
destroyed, we can assign it a new content. // We will, in this case,
assign an empty value to the vector, // making it effectively empty vec
= \{\}; // Since the vector as gained a determinate value, we can use it
normally. vec.push\_back(42); // Send the vector by move again.
consumingFunction(std::move(vec)); \} Section 106.5: Move assignment
Similarly to how we can assign a value to an object with an lvalue
reference, copying it, we can also move the values from an object to
another without constructing a new one. We call this move assignment. We
move the values from one object to another existing object. For this,
we will have to overload operator =, not so that it takes an lvalue
reference, like in copy assignment, but so that it takes an rvalue
reference. class A \{ int a; A\& operator= (A\&\& other) \{
this-\textgreater a = other.a; other.a = 0; return }this; \} \}; This is
the typical syntax to dene move assignment. We overload operator = so
that we can feed it an rvalue reference and it can assign it to another
object. A a; a.a = 1; A b; b = std::move(a); //calling A\& operator=
(A\&\& other) std::cout \textless\textless{} a.a \textless\textless{}
std::endl; //0 std::cout \textless\textless{} b.a \textless\textless{}
std::endl; //1 Thus, we can move assign an object to another one.
Section 106.6: Using move semantics on containers You can move a
container instead of copying it: void print(const std::vector\& vec) \{
for (auto\&\& val : vec) \{ std::cout \textless\textless{} val
\textless\textless{} ``,''; \} std::cout \textless\textless{} std::endl;
\} int main() \{ // initialize vec1 with 1, 2, 3, 4 and vec2 as an empty
vector std::vector vec1\{1, 2, 3, 4\}; std::vector vec2; // The
following line will print 1, 2, 3, 4 print(vec1); // The following line
will print a new line print(vec2); // The vector vec2 is assigned with
move assingment. // This will ``steal'' the value of vec1 without
copying it. vec2 = std::move(vec1); // Here the vec1 object is in an
indeterminate state, but still valid. // The object vec1 is not
destroyed, // but there's is no guarantees about what it contains. //
The following line will print 1, 2, 3, 4 print(vec2); \}
